%=====================================================================
%  Method section --- full standalone LaTeX (compiles as-is).
%
%  This document uses the standard `article' class in two-column mode so
%  it builds anywhere with no external style file.  For the real NeurIPS
%  submission, delete the \documentclass line and the \geometry line
%  below and use instead:
%       \documentclass{article}
%       \usepackage[final]{neurips_2024}
%  then keep everything from \begin{document} onward unchanged.
%=====================================================================
\documentclass[10pt,twocolumn]{article}

\usepackage[letterpaper,margin=0.85in,columnsep=0.28in]{geometry}
\usepackage{amsmath,amssymb}
\usepackage{booktabs}
\usepackage{microtype}
\usepackage[hidelinks]{hyperref}

\DeclareMathOperator*{\argmin}{arg\,min}
\DeclareMathOperator{\tr}{tr}
\newcommand{\R}{\mathrm{R}}   % Refusal
\newcommand{\J}{\mathrm{J}}   % Jailbreak
\newcommand{\B}{\mathrm{B}}   % Benign

\begin{document}

% Full-width title + abstract spanning both columns (standalone-safe).
% With the neurips class, replace this \twocolumn[...] block with the usual
% \maketitle and \begin{abstract}...\end{abstract}.
\twocolumn[{%
  \begin{center}
    {\LARGE\bfseries Coordinate Geometry of Activation-Space
      Decision Boundaries for Jailbreak Detection\par}
    \vspace{0.35em}
    {\large\itshape Method\par}
    \vspace{0.6em}
    {\normalsize Anonymous Author(s)\par}
  \end{center}
  \vspace{0.4em}
  \begin{center}
  \begin{minipage}{0.86\textwidth}
    \small\noindent\textbf{Abstract.}\ \
    We treat the \emph{placement} of the refusal/jailbreak decision boundary in a language
    model's residual stream as the object of measurement, and study it in three experiments.
    Experiment~1 collects behaviorally-labeled last-token activations and unwraps their geometry
    into a spherical chart, yielding one site per class. Experiment~2 fits a family of
    covariance-aware partitions over those sites---ranging from scalar-weighted power diagrams to
    the full-covariance (QDA) limit---and calibrates a geometric margin. Experiment~3 benchmarks
    each boundary against a linear probe, measures how it transfers to unseen attack families, and
    diagnoses the class geometry. Throughout, all quantitative geometry is computed in a reduced
    Euclidean or spherical space; two-dimensional embeddings are used only for display.
  \end{minipage}
  \end{center}
  \vspace{1.2em}
}]

\section{Method}
\label{sec:method}

Consider an instruction-tuned language model whose residual stream has width $d$. For a prompt
$p$ we read off a single hidden state $x \in \mathbb{R}^{d}$ at each layer, and we label that
state by \emph{what the model did}: Refusal ($\R$), Jailbreak ($\J$), or Benign ($\B$). Within
whatever coordinate system we are currently using, each class has an empirical mean $\mu_c$ and
covariance $\Sigma_c$.

The detection question---is this activation a jailbreak?---is really a question about a
surface. A linear probe draws that surface as a flat hyperplane, and the direction of the
hyperplane is fixed by the discriminative objective it was trained on. A class-conditional
method instead derives the surface from the class statistics $\{\mu_c, \Sigma_c\}$. These two
surfaces rarely coincide, and the quantity we care about is the gap between them: how far the
probe's hyperplane sits from the covariance-correct boundary, and how that gap grows as the
class covariances diverge. The three experiments below are built to make that gap measurable.

\subsection{Experiment 1: unwrapping the geometry}
\label{sec:e1}

\paragraph{Data and labels.}
We draw prompts in equal quotas from four safety corpora---SORRY-Bench, WildJailbreak,
ToxicChat, and Aegis~2.0---so that no single jailbreak family dominates the sample. Crucially,
the class of an example is decided by \emph{behavior}, not by its source label: $\R$ is a
harmful prompt the model refused, $\J$ a harmful prompt it complied with, and $\B$ a harmless
prompt it complied with. A harmless prompt that the model refuses is over-refusal; we hold it
out of the training partition and use it only as a stress test in Experiment~3. Refusal is
detected from the first few generated tokens, which lets us keep generation short and batched.
Because a compliant model refuses only a minority of harmful prompts, we over-sample the
harmful side, label everything behaviorally, and then downsample to a common class
size---canonically $n = 500$ per class.

\paragraph{Representation.}
Our reference model is Qwen2.5-3B-Instruct (36 layers, loaded in 4-bit NF4). At every even
layer $L \in \{2, 4, \dots, 36\}$ we store the residual state of the \emph{last} token. This
choice matters: mean-pooling the sequence washes out the very signal we rely on, driving the
inter-class covariance similarity to $0.95$--$0.999$, whereas the last-token state keeps the
classes distinctly shaped. Estimating a full-width covariance from a few hundred points is
ill-posed, so we first reduce each layer to $k$ principal components (canonically $k = 10$).
This reduced Euclidean space is our first working coordinate system.

\paragraph{Spherical unwrapping.}
Our central hypothesis is that the residual stream lies close to a hypersphere, and that the
mismatch between class covariances is largely \emph{radial}. If that is true, then discarding
the radius should equalize the remaining (tangential) covariances---and a flat boundary, wrong
in the ambient space, becomes correct once the sphere is unwrapped. We test this with a custom
implementation of Principal Nested Spheres (PNS)~\citep{jung2012pns}. After projecting the
points onto the unit sphere $S^{k-1}$, we peel off nested subspheres one at a time; each
subsphere is an axis $v$ and a radius-angle $r$ chosen to hug the data,
\begin{equation}
  (v, r) \;=\; \argmin_{v, r}\ \sum_i \big(\angle(x_i, v) - r\big)^2 ,
  \label{eq:pns}
\end{equation}
where $\angle(x, v) = \arccos\langle x, v\rangle$ is the geodesic angle. Stacking the signed
residual angles from every peeled subsphere gives a Euclidean \emph{PNS-score} representation,
our second working coordinate system. In either space we then summarize each class by its
\emph{site}: the mean $\mu_c$ and a shrunk covariance $\Sigma_c$. A per-layer figure that places
the raw sphere beside its unwrapped version confirms the picture the boundary will exploit:
Refusal collapses to a tight cap, while Jailbreak fans out.

\subsection{Experiment 2: the boundary and its margin}
\label{sec:e2}

Given sites $\mu_i$ with scalar weights $w_i$, the power (Laguerre) distance and the boundary
between sites $i$ and $j$ are
\begin{align}
  \pi_i(x) &= \|x - \mu_i\|^2 - w_i , \label{eq:power}\\
  2(\mu_j - \mu_i)\!\cdot\! x &= \|\mu_j\|^2 - \|\mu_i\|^2 - w_j + w_i . \label{eq:hyper}
\end{align}
The boundary in \eqref{eq:hyper} is a hyperplane whose normal is $(\mu_j - \mu_i)$; the weight
$w_i$ can only slide it back and forth along that normal. We set $w_c = \tr(\Sigma_c)$, which
pushes the fence away from a diffuse class and toward a tight one. Note what a \emph{single}
site per class cannot do: the $\|x\|^2$ terms in \eqref{eq:power} cancel, the boundary stays
exactly flat, and so it can never bend to match the curved Bayes boundary that unequal
covariances demand. To recover that curvature while staying within the power-diagram family, we
split each class into $m$ sub-sites by $k$-means (canonically $m = 12$), give each its own
scalar weight, and score a class by its nearest sub-site,
\begin{equation}
  s_c(x) = \min_{s \in c}\big( \|x - \mu_s\|^2 - w_s \big).
  \label{eq:multisite}
\end{equation}
The resulting inter-class boundary is piecewise linear: a union of hyperplane facets that
approximates a curved surface while remaining a genuine Laguerre partition with locally convex
cells.

A tempting alternative is to fold the whole covariance into the distance,
\begin{equation}
  P_c(x) = (x - \mu_c)^{\!\top} \Sigma_c^{-1} (x - \mu_c) + \log\det \Sigma_c .
  \label{eq:qda}
\end{equation}
We include \eqref{eq:qda} in our comparison, but we are careful to call it what it is: this is
exactly quadratic discriminant analysis (QDA). Its boundary is the quadric $P_A = P_B$, and its
cells need not be convex. The distinction matters because the Alexandrov--Minkowski theorem
certifies a unique, mass-correct optimal-transport partition only for the \emph{scalar-weight}
diagrams \eqref{eq:power}--\eqref{eq:multisite}; it says nothing about the full-covariance form.
All covariances are estimated with Ledoit--Wolf shrinkage so that $\Sigma^{-1}$ and $\log\det
\Sigma$ remain well-defined at small $n$. Table~\ref{tab:family} lists the family.

\emph{Why a nonlinear chart is necessary.} Two symmetric positive-definite covariances can
always be simultaneously diagonalized, but they can be simultaneously \emph{whitened} only when
$\Sigma_i \propto \Sigma_j$. A linear change of coordinates can therefore shrink the mismatch to
a spectrum of eigenvalue ratios, but it can never set that spectrum to one. If covariance
mismatch is what displaces the boundary, then removing it requires the nonlinear chart of
Experiment~1---not a cleverer linear reparameterization.

\begin{table}[t]
  \centering
  \small
  \caption{The discriminant family compared in Experiments~2--3. Only the scalar-weight forms
  are true power diagrams and carry the optimal-transport guarantee.}
  \label{tab:family}
  \begin{tabular}{@{}lll@{}}
    \toprule
    Discriminant & Boundary & Power diagram / OT \\
    \midrule
    LDA (pooled $\Sigma$)  & linear            & no (covariance-blind ref.) \\
    Power, single-site     & linear            & yes \\
    Power, multi-site      & piecewise-linear  & yes \\
    Full-covariance (QDA)  & quadric           & no (this is QDA) \\
    \bottomrule
  \end{tabular}
\end{table}

\paragraph{Margin.}
Sorting a point's class scores from smallest to largest, we assign it to the smallest and read
its confidence off the gap to the runner-up,
\begin{equation}
  \gamma(x) = s_{(2)}(x) - s_{(1)}(x).
  \label{eq:margin}
\end{equation}
A large $\gamma$ means the point sits deep inside its cell; a small one means it sits on the
boundary. For a power diagram this is a signed power-distance gap; for QDA it is a
log-likelihood-ratio gap. Because these margins are unbounded and vary in scale across layers and
methods, we rank-normalize each to a confidence in $[0,1]$ within its run, bin the confidence,
and compare bin confidence to bin accuracy through the expected calibration error,
\begin{equation}
  \mathrm{ECE} = \sum_b \frac{n_b}{n}\,\big|\operatorname{acc}(b) - \operatorname{conf}(b)\big|.
  \label{eq:ece}
\end{equation}
We do not expect the in-distribution ECE to separate the methods; the property we are after is
whether the margin stays honest \emph{under distribution shift}.

\paragraph{How far is the boundary from correct?}
To say a boundary is well placed we need a reference to place it against. We fit a per-class
Gaussian mixture ($g$ components, canonically $g = 2$) as that reference, and take its
class-likelihood crossing as the reference boundary. Walking along the segment that joins two
class means, we find where the method flips its preference ($t_m$) and where the reference flips
its own ($t_r$), and report the RMS displacement between the two crossings, scaled by the
distance between the means,
\begin{equation}
  \varepsilon_b = \sqrt{\ \sum_{(a,b)} \Big[(t_m - t_r)\,\|\mu_a - \mu_b\|\Big]^2 }.
  \label{eq:rms}
\end{equation}
Alongside $\varepsilon_b$ we report a scale-free companion---the fraction of points that the
method and the reference assign to different cells---since it, unlike $\varepsilon_b$, is
comparable across coordinate systems. One case is deliberately excluded: when the method is QDA
and the reference is a single Gaussian per class ($g = 1$), the two are the same object and
$\varepsilon_b \approx 0$ for trivial reasons. Our pipeline flags exactly this configuration and
never counts it as a win.

\subsection{Experiment 3: benchmarks and diagnostics}
\label{sec:e3}

We turn each discriminant into a binary detector that emits one score for the positive class, and
compare them all under a single ROC-AUC against two baselines: a logistic-regression
\emph{linear probe}, the standard activation detector, and a \emph{plain Voronoi} classifier
(nearest class mean, no covariance weight). The primary task is Jailbreak versus Benign---the
hard pair, since both are compliance and overlap heavily---scored by stratified 5-fold AUC at
each layer. Plotting AUC against depth gives the paper's decisive test: does the geometry
method's advantage over the flat probe \emph{widen} in the late layers, where the boundary is
most misplaced?

Two further probes round out the picture. For \emph{cross-family} transfer we leave one source
corpus out, train on the remaining three, and report the drop from in-distribution to held-out
AUC as a direct measure of how each detector survives an unseen attack style. For the
\emph{fractal diagnostic} we measure each class's effective number of spread directions with the
covariance participation ratio,
\begin{equation}
  D = \frac{\tr(\Sigma)^2}{\|\Sigma\|_F^2}
    = \frac{\big(\sum_i \lambda_i\big)^2}{\sum_i \lambda_i^2},
  \label{eq:frac}
\end{equation}
and test the prediction that the diffuse class occupies many more directions than the tight one,
$D(\J) \gg D(\R)$. We use \eqref{eq:frac} rather than box-counting, which is unstable in high
dimension with few points.

Every boundary in Experiments~2--3 is fit in both the raw-PCA space and the PNS space, so that
the coordinate-dependence at the heart of the paper is measured rather than assumed. Layer
stability---the depth at which AUC levels off---is read directly from the AUC-versus-depth curve
and reported as an early-exit safety horizon. Finally, every source of randomness (PCA, $k$-means,
the mixtures, and all data splits) uses a fixed seed, and collection (Experiment~1) is kept
separate from analysis (Experiments~2--3) so that every geometric result can be reproduced on CPU
from the saved activations.

\paragraph{Scope of the construction.}
We keep two claims apart. First, the optimal-transport and convex-cell guarantees hold for the
scalar-weight power diagrams in \eqref{eq:power}--\eqref{eq:multisite}, but not for the
full-covariance form in \eqref{eq:qda}, which is QDA. Second, whether a covariance-aware boundary
actually \emph{beats} a linear probe---and in which coordinate system---is an empirical matter
that Experiment~3 decides; it is not guaranteed by the construction. We therefore report the
AUC-versus-depth and cross-family curves in both spaces and let the measurements settle the
framing.

%---------------------------------------------------------------------
\begin{thebibliography}{1}
\small
\bibitem{jung2012pns}
S.~Jung, I.~L. Dryden, and J.~S. Marron.
\newblock Analysis of principal nested spheres.
\newblock \emph{Biometrika}, 99(3):551--568, 2012.
\end{thebibliography}

\end{document}
